% Invented sample: every name, number and reference is made up.
\chapter{Sizing the return desks at a tool library}
\label{ch:desks}

\section{Introduction}
\label{sec:intro}

The Millbrook Tool Library lends tools to its members for a small yearly fee.
Members borrow drills, ladders and garden tools, and they bring them back to a return desk, where a volunteer checks each tool for damage before it goes back on the shelf.
On busy days a queue forms at the desk, and some members leave without returning their tools.
Most tool libraries in the region report the same problem.

This chapter applies the queueing model from Chapter~\ref{ch:background} to the return desk at Millbrook.
In this chapter we will first set up the model, then derive the waiting times, then describe a staffing rule and finally compare the model with the desk logs.

The key quantity is the utilisation $\rho = \lambda/(N\mu)$, which says how busy the desks are.
When $\rho$ is close to one, even small changes in demand lead to long queues.

The library opens on weekday afternoons from 14:00 to 18:00 and on Saturdays from 9:00 to 13:00.
Returns peak on Saturday mornings, when members bring back the tools they borrowed for weekend projects.

\section{The model}
\label{sec:model}

We assume that returns arrive as a Poisson process with rate $\lambda$ returns per hour, and that each volunteer serves returns at rate $\mu$ returns per hour, with exponential service times.
The library opens $N$ desks, each staffed by one volunteer.
These are the assumptions of Chapter~\ref{ch:background}, and they give the standard M/M/$N$ queue \cite{example2021}.
Exponential service times suit this setting, because a check can take a few seconds for a clean trowel or several minutes for a hedge trimmer with a damaged cable \cite{sample2019}.

The desk logs give the parameter values.
A volunteer takes 3 minutes on average to check a returned tool, so $\mu = 20$ returns per hour.
The library has space for 2 desks, so $N$ is either one or two.

With these values the utilisation is
\begin{equation}
  \rho = \frac{\lambda}{N\mu},
  \label{eq:util}
\end{equation}
as in Equation~\eqref{eq:rho}.
A borrower has to wait whenever all $N$ desks are busy.
For one desk, this happens with probability $P_w = \rho$.
For two desks, the Erlang C formula gives
\begin{equation}
  P_w = \frac{2\rho^2}{1+\rho}.
  \label{eq:erlang2}
\end{equation}

The time in system $W$ grows quickly as $\rho$ approaches one, which is why the choice between one and two desks matters.
With the values above, the model gave a mean wait of 2.0 minutes on weekday afternoons, which matched the desk logs closely.
Section~\ref{sec:waiting} compares the model with the desk logs.

\section{Waiting times}
\label{sec:waiting}

Two measures of delay matter to a borrower.
The waiting time $W_q$ is the time a borrower spends in the queue before a volunteer starts to check the tool.
The time in system $W$ also includes the check itself, so
\begin{equation}
  W = W_q + \frac{1}{\mu}.
  \label{eq:w}
\end{equation}
The library's target concerns the waiting time, not the time in system, because members do not mind the check itself.
The target is a mean waiting time of less than 3 minutes.
Earlier studies of public service desks use the same target.

It is worth noting that, for a queue with $N$ desks, the mean waiting time can clearly be obtained from the probability of waiting in Equation~\eqref{eq:erlang2} by dividing it by the spare service capacity $N\mu - \lambda$, as shown in the very well known result \cite{example2021}, which gives
\begin{equation}
  W_q = \frac{P_w}{N\mu - \lambda},
  \label{eq:wq}
\end{equation}
and this is the formula that I use throughout the rest of the chapter, including in Section~\ref{sec:logs}.

Little's law links the mean number of borrowers at the desk area to the time in system:
\begin{equation}
  N = \lambda W.
  \label{eq:little}
\end{equation}
For one desk and $\lambda = 8$ returns per hour, Equation~\eqref{eq:wq} gives $W_q = 0.4/12$ hours, or 2.0 minutes, and the desk area holds about $N = 0.67$ borrowers on average.
At this load the desk is half asleep.
This result shows the wider value of queueing models for community services.

\section{A staffing rule}
\label{sec:rule}

% keep light: adapted from the published paper writer2025
This section is adapted from \cite{writer2025}, which I wrote with my supervisor.
I built the model and ran the analysis, and my supervisor suggested the form of the rule.

We propose a simple staffing rule.
Open one desk while returns stay below 10 per hour, and open a second desk when they rise above that level.
The rule follows from Equation~\eqref{eq:wq}, since with one desk the waiting time stays under 3 minutes only while $\lambda < 10$, and it is clearly easy for volunteers to apply because the desk logs already record the number of returns in each hour of the day.

With two desks, the same target holds up to about 28 returns per hour, which is well above the busiest hour in the logs.
% end keep

\section{Checking the model against the desk logs}
\label{sec:logs}

The library kept desk logs for 16 weeks.
Each log records the time each borrower joined the queue and the time the check began, so the observed waiting time can be compared with $W_q$ from our model.
The logs hold 3,296 returns. Of these, 41 had no start time and are left out.
On weekday afternoons returns ran at 4 to 9 per hour; on Saturday mornings they ran at 16--22 per hour.

The library opens on weekday afternoons from 14:00 to 18:00 and on Saturdays from 9:00 to 13:00.
Returns peak on Saturday mornings, when members bring back the tools they borrowed for weekend projects.

Table~\ref{tab:logs} compares the observed and predicted waiting times.
For the Saturday morning average with two desks, the predicted waiting time is
\begin{equation}
  W_q = \frac{P_w}{2\mu - \Lambda} = \frac{0.279}{40 - 18}\ \text{hours},
  \label{eq:saturday}
\end{equation}
or 0.76 minutes, against 0.9 minutes in the logs.

\begin{table}[h]
  \centering
  \begin{tabular}{lccc}
    Period & Desks & Predicted $W_q$ (min) & Observed $W_q$ (min) \\
    \hline
    Weekday afternoon ($\lambda = 8$) & 1 & 2.0 & 2.1 \\
    Saturday morning ($\lambda = 18$) & 2 & 0.76 & 0.9 \\
    Saturday peak hour ($\lambda = 22$) & 2 & 1.3 & 1.7 \\
  \end{tabular}
  \caption{Predicted and observed waiting times at the return desk.}
  \label{tab:logs}
\end{table}

The model slightly underestimates the waiting time at the Saturday peak.
The time in system is less affected, because the check itself takes the same 3 minutes whatever the length of the queue.
The largest gap is at the busiest hour, where the logs show 1.7 minutes against 1.3 predicted.
Service times at the peak are less regular than the exponential assumption allows, because volunteers check heavy tools more slowly when the queue is long.

The rule is sensitive to the service rate.
If a check took 4 minutes on average, so that $\mu = 15$ returns per hour, one desk would meet the target only below about 6 returns per hour.

\section{Conclusion}
\label{sec:conclusion}

One desk keeps the mean waiting time under 3 minutes as long as returns stay below 10 per hour.
Above that, a second desk is needed, and two desks keep the waiting time under 3 minutes up to about 28 returns per hour.
The desk logs agree with the model to within half a minute.
For Millbrook, this means opening a second desk on Saturday mornings only.
The next chapter extends the model to borrowers who leave the queue before they are served.
